% Esquisse d'un programme (cote 119, copy 1): archive cote 119, site page 30 (PDF page 31), typescript page 28.
% typescript. Draft compared: EsquisseFr.pdf (text layer):1019-1070.
% AI-transcribed by Claude Opus 5.5 (claude-opus-5-5, adaptive thinking, effort high), Message Batches API, batch pilot (msgbatch_0189Nhf1xB615YWnztubtLfc), 2026-11-12; not yet checked by eye.
\scanpage{119/30}
\mspage{28}

les classes d'isotopie (dans un sens dicté sans ambiguité par le
contexte) d'isomorphismes (voire même, de morphismes plus généraux
que des isomorphismes). Cette idée, qui est reprise avec succès
dans la thèse de Yves Ladegaillerie notamment (cf. début du par. 3),
m'a servi de façon systématique dans toutes mes réflexions ulté-%
rieures de topologie combinatoire, quand il s'est agi de formuler
avec précision des théorèmes de traduction de situations topolo-%
giques, en termes de situations combinatoires. Dans la situation
actuelle, mon espoir était d'arriver à formuler (et à prouver !)
un théorème d'équivalence entre deux catégories isotopiques conve-%
nables, l'une étant la catégorie des ``paires admissibles'' $(X,Y)$,
l'autre celle des ``triples admissibles'' $(U,Y,f)$, où $Y$ est une
variété, $U$ une variété à bord, et $f : \partial U \longrightarrow Y$ une fibration
propre et lisse. De plus, bien sûr, j'espérais qu'un tel énoncé,
modulo un travail de nature essentiellement algébrique, s'éten-%
drait de lui-même en un énoncé plus sophistiqué, s'appliquant aux
structures stratifiées générales.

Très vite il apparaissait qu'il ne pouvait être ques-%
tion d'obtenir un énoncé aussi ambitieux dans le contexte des
espaces topologiques, à cause des sempiternels phénomènes de
``sauvagerie''. Déjà quand $X$ lui-même est une variété et $Y$ ré-%
duit à un point, on se bute à la difficulté que le cône sur un
espace compact $Z$ peut être une variété en son sommet, sans que
$Z$ soit homéomorphe à une sphère, ni même soit une variété. Il
était clair également que les contextes de structures plus rigides
qui existaient à l'époque, tel le contexte ``linéaire par morceaux'',
étaient également inadéquats - une des raisons rédhibitoires com-%
munes étant qu'ils ne permettaient pas, pour une paire $(U,S)$
d'un ``espace'' $U$ et d'un sous-espace fermé $S$, et une application
de recollement $f : S \longrightarrow T$, de construire la somme amalgamée cor-%
respondante. C'est quelques années plus tard que j'étais informé
de la théorie de Hironaka des ensembles qu'il appelle, je crois,
``semi-analytiques'' (réels), qui satisfont à certaines des condi-%
tions de stabilité essentielles (sans doute même à toutes) néces-

\margin{30}
